% ============================================================
%  R. Nielsen, "Et Synspunkt for Darwinismen"
%  For Idé og Virkelighed, 1873, pp. 449–459
%  TRANSCRIPTION (Danish) — from Google Books scan
% ============================================================
\documentclass[12pt,a4paper]{article}
\usepackage[utf8]{inputenc}
\usepackage[T1]{fontenc}
\usepackage[danish]{babel}
\usepackage{microtype}
\usepackage{setspace}
\usepackage[margin=1.2in]{geometry}
\usepackage{fancyhdr}
\usepackage{hyperref}

\hypersetup{
  pdftitle={Et Synspunkt for Darwinismen},
  pdfauthor={R. Nielsen},
  hidelinks
}

\pagestyle{fancy}
\fancyhf{}
\rhead{Nielsen, \emph{Et Synspunkt for Darwinismen} (1873)}
\cfoot{\thepage}

\setstretch{1.3}

\title{\textbf{Et Synspunkt for Darwinismen}}
\author{R.\ Nielsen}
\date{\emph{For Idé og Virkelighed}, 1873, s.\ 449--459}

\begin{document}
\maketitle
\thispagestyle{fancy}

\bigskip

\noindent Darwin er ikke Naturfilosof, men --- i udmærket Forstand ---
Naturforsker. Denne simple Bemærkning, der ikke antyder nogen ny
Opdagelse, men blot udsiger, hvad Enhver, der befatter sig med at
gjøre Anmeldelser af Darwinshypothesen, kan sige sig selv, er efter
vort Skjøn vel skikket til at fjerne en Del af de Misforstaaelser,
hvormed en gjærende Opinion fra mangfoldige Sider synes at ville
belemre et i sin Retning storartet Forsøg.

De Synspunkter, hvorom Opinionspolemiken dreier sig, ere:

\begin{enumerate}
  \item Organismernes fælles Afstamning, særlig Menneskets Oprindelse.
  \item Fremkomsten af de første organiske Væsener paa Jorden,
        hvorvidt den skyldes en Skabelse eller blot en eiendommelig
        Samvirken af almindelige Naturkræfter.
  \item Darwinismens Karakter i det Hele, om den virkelig kan bestaae
        med Menneskets Værdighed, eller maa siges at indeholde
        aandsfornegtende, for Religion og Sædelighed farlige Tendenser.
\end{enumerate}

\noindent Til Opklaring af disse Hovedpunkter have vi, hvad navnlig
det sidste angaaer, aldeles ingen Veiledning i Darwins rent personlige
Overbevisning, der ikke er os nærmere bekjendt og ei heller vedkommer
Sagen; vi holde os udelukkende til Theorien, saaledes som den
foreligger.

% [Side 450]

\medskip

\noindent 1) »Hvorfra ere disse Legioner af levende Væsener, disse næsten
talløse Arter og Slægter af Planter og Dyr, som findes i Naturen?«
Darwin svarer: »de nedstamme alle fra faa, simpelt organiserede Former;
de udgjøre eet sammenhængende System, hvis enkelte Led, tiltagende i
Fylde og Fuldkommenhed, have i Tidernes Løb arbeidet sig frem under en
vedvarende Kamp for Livet.«

Det vil nu let kunne sees, at det gjør en stor Forskjel, om vi opfatte
den saaledes udtalte Anskuelse fra Erfaringslærens eller fra
Naturfilosofiens Standpunkt. Anlægger man nu en naturfilosofisk
Maalestok, saa spørges der med Grund om det Princip, hvoraf Theorien i
sin Helhed lader sig aflede.

I denne Henseende have ikke blot Filosoferne, men selv Naturforskere,
som f.\ Ex.\ Th.\ L.\ Bischoff, fundet sig utilfredsstillede. »Den
Darwinske Lære«, siger Bischoff, »er ingen Theori, den hviler ikke paa
et Princip, som fordi det er afledet som et almengyldig rigtigt, ogsaa
maa af Alle erkjendes som et almengyldig sandt. Den indeholder meget
vigtige og sande Tanker og Kjendsgjerninger; men disse kan man kun
anerkjende for rigtige, for saa vidt de for den objektive Erfaring og
Iagttagelse bekræfte sig som rigtige. Men hvor denne Eftervisning ikke
lykkes eller ikke er lykkedes, ville vi ikke kunne tillægge den
Darwinske Theori nogen Betydning.«

% [Side 451]

\medskip

\noindent Froschhammer\footnote{Om Darwinismen, oversat af H. Steen.
Kbhvn. 1873.}, der bedømmer Darwinismen udelukkende fra et filosofisk
kritisk Synspunkt, lægger blandt Andet Vægt paa den Uklarhed, der
hviler over Darwins bekjendte Slagord: »den naturlige Udvælgelse«
(Kvalitetsvalget). Da en Del af hans øvrige Indvendinger forekomme os
at være af temmelig tvivlsom Betydning, skulle vi særlig fremhæve
denne, der, i naturfilosofisk Forstand, vel at mærke, ganske vist er
væsenlig begrundet.

Fr.\ vælger som Exempel Øiets Oprindelse og Udvikling. »Da vi«, siger
han, »maae tænke os de primitive Dyreformer, simple og ufuldkomne som
de ere, uden Øine, ligesom der jo endnu existerer Dyr uden Øine, saa
opstaaer paany Hovedspørgsmaalet, hvordan da fra første Begyndelse af
Øinene opstode eller kunde opstaae. De maatte enten opstaae ved
Tilfælde, eller ved en uforklarlig, ubegribelig generatio æqvivoca,
eller udtrykkelig ved en ny Skabelsesakt. I alle Tilfælde kunde de
ikke opstaae ved en naturlig Udvælgelse, da denne efter sit Begreb kun
formaaer at modificere eller egenlig kun at vedligeholde et allerede
Givet, ikke derimod at skabe et Nyt, hidtil ikke Givet.«

Det fuldkomne Øie sammenlignes med Teleskopet og den naturlige
Udvælgelses Virksomhed for at fuldkommengjøre Øiet med den menneskelige
Intelligenses Anstrængelser for at forbedre hint Instrument; »men« ---
tilføier den filosofiske Kritiker --- »sikkert med Urette. Den bevidstløse
Natur er lige saa lidt i Stand til at efterligne eller udøve Optikernes
planmæssige Virksomhed, som den

% [Side 452]

er i Stand til at efterligne eller erstatte Konstnerens (Malerens eller
endog blot Uhrmagerens) Virksomhed. Darwin geraader paa dette Sted,
for ikke at gaae i Staa med sin Forklaring af de fuldkomneste Øines
Oprindelse, ind paa en formelig Personifikation af den naturlige
Udvælgelse. Den »naturlige Udvælgelse« skal »nøie iagttage« og
»omhyggelig udvælge« og med aldrig feilende Takt vide at udfinde
enhver Forbedring, der kan befordre den videre Fuldkommengjørelse«.
»Var dette at forstaae ligefrem, saa var af Darwin selv en theologisk
Magt indført i Naturen, som vilde gjøre alle hans øvrige
Forklaringsforsøg overflødige?« »men tillige en Egenskab tillagt den
naturlige Udvælgelse, som ganske staaer i Strid med dens øvrige Væsen.
Men ere de nævnte Udtryk uegentlig at forstaae, saa er her i Ord givet
eller fingeret en Forklaring, idet kun Noget figurlig paastaaes, som
ikke i Realiteten kan finde Sted\ldots\ Den naturlige Udvælgelse
kan\ldots\ ikke tilstræbe fuldkomnere Øine, men kan kun vedligeholde
og benytte dem, naar de engang paa en eller anden Maade ere blevne
til. Og her synes da altsaa paa slaaende Maade det Tilfælde givet,
hvorom Darwin selv tilstaaer, at det kunde omstyrte hans Theori.«
Lignende Indsigelser kunne reises imod den Maade, hvorpaa Menneskets
Oprindelse (i det udførlige Værk: \emph{The Descent of Man}) bliver
forklaret. For at forberede Læseren paa den kommende Udvikling, bliver
strax i Begyndelsen af første Del et Menneskefoster sammenstillet med
et Hundefoster; fra Ligheden i det Ydre sluttes der til Lighed i
Væsen, en Slutningsmaade, der er gjennemgaaende og karakteristisk for
hele den følgende Bevisførelse. Vil Nogen kaste sig med kritisk Iver
over dette Punkt, vil

% [Side 453]

det neppe være ham vanskeligt at opdage og blotte de principielle
Svagheder, naar Sagen, vel at mærke, behandles filosofisk.

Men stille vi os paa Forskerens Standpunkt, faaer det Hele et andet
Udseende. Naturforskeren har ikke den Opgave at konstruere de levende
Væseners typiske Eiendommeligheder og stedfindende Forskjelligheder ud
af et almindeligt Princip; tvertimod! han forefinder de forskjellige
Former givne i Naturen. Hans første Spørgsmaal er ikke et
Nødvendigheds-, men et Mulighedsspørgsmaal. Skulle de mange forskjellige
Arter og Slægter antages at staae uafhængige af hverandre, da maa der
forudsættes ligesaa mange Skabelsesmirakler som der er Arter; men at
bygge paa Mirakler er at opgive al Forskning; hver Gang man er i
Forlegenhed med at udfinde en naturlig Grund til en naturlig
Kjendsgjerning, kan man jo paa den allerbekvemmeste Maade bryde af ved
at sige: her er skeet et Mirakel. »Efter den almindelige
Betragtningsmaade, den, at hvert enkelt Væsen er Resultat af en egen
Skabelsesakt, kunne vi kun sige, at saaledes er det, at det har
behaget Skaberen at indrette alle Dyr og Planter i hver stor Klasse
efter en ensformig Plan; men dette er ikke nogen videnskabelig
Forklaring.«

Disse Ord af Darwin angive tydelig nok, hvad det er, hvorpaa han
sigter. Darwins Lære optræder ikke som nogen fuldfærdig Forklaring og
Begrundelse af de organiske Typers stedfindende Eiendommelighed, men
som en Hypothese. Denne Hypothese har den store videnskabelige
Betydning, at den leder Forskningen ind i nye Spor og angiver en ny
Methode. At opstille Hypotheser er en let Sag; men at angive nye og
frugtbare Methoder er forbeholdt Videnska-

% [Side 454]

bens Udmærkede. Hvad Darwin lærer os om Organernes Udvikling og
Omdannelse i Kampen for Livet, om Kvalitetsvalg, Arvelighed,
Varieringslove o.\ s.\ v., ere afgjørende Bestemmelser ved den
Retning, hvori hans Forskninger gaae, altsaa ved hans Methode.
Istedenfor at paaprakke ham et naturfilosofisk Princip, som han ikke
har, gjorde man bedre i at fæste sin hele Opmærksomhed paa Methodens
Eiendommelighed, thi Forskerens Princip er hans Methode.

Darwin er ikke gjendreven, fordi han endnu ikke har kunnet paavise,
hvorledes et Bløddyr er gaaet over til at blive enten et Leddyr eller
et Hvirveldyr og endnu mindre, hvorledes en Abe med filosofisk
Nødvendighed er blevet til et Menneske; han sigter, saavidt vi have
forstaaet hans Methode, ikke paa de tilblevne Formers Nødvendighed,
men paa de gradvise Overganges Mulighed og derved paa Muligheden af at
naae til større og større Indsigt i de organiske Formers indbyrdes
naturlige Sammenhæng. Hvad de enkelte Mellemled angaaer, er der Ingen,
der stærkere end Darwin har indskærpet, hvor umaadelig stor vor
Uvidenhed er.

% [Side 455] (numbered as 2nd main question in original)

\medskip

\noindent 2) Hvorledes ere de første organiske Væsener fremkomne paa
Jorden? Dette Spørgsmaal kan ingen Naturforsker besvare, og det
forekommer os, at Darwin paa den værdigste Maade har hævdet sin
Videnskabs Standpunkt ved slet ikke at give sig Mine af at ville
besvare det. »Der er«, siger han, »Storhed i det Syn paa Livet, at
det med dets forskjellige Kræfter, af Skaberen oprindelig er blevet
nogle faa eller en enkelt Form indblæst, og at, medens denne vor Klode
har rullet rundt efter Tyngdens bestemte Lov, have utallige Former,
høist skjønne og høist vidunderlige, fra en simpel Begyndelse udviklet
sig og udvikle sig endnu.«

I Sandhed et Sprog, der er en Naturforsker værdigt! Skal
Darwinshypothesen absolut skrues ind i et naturfilosofisk System og
underkaste sig en »dialektisk Kritik«, da er det unegtelig sandt, at
Forskeren maa komme i Forlegenhed. Har man først indrømmet eet Mirakel,
hvorfor saa ikke indrømme flere? Ere de første organiske Væsener
skabte, hvorfor maae de da ikke alle være skabte? Det er en Kritik,
som Enhver let kan overkomme, og det er derfor ikke at undres over, at
den har ladet sig høre fra saa mange Sider. Men hvad er nu egenlig
Forskerens Tanke? Han seer i Livet en Virksomhed, som vel overalt
falder ind under almengjældende Naturlove, men dog en Virksomhed, som
ikke lader sig aflede eller forklare af elementære Kræfters
Vexelvirkning. Livets Oprindelse er en Gaade, og Gaaden antydes ved,
at »Livet med dets forskjellige Kræfter af Skaberen oprindelig er
blevet nogle faa eller en enkelt Form indblæst«. Ved denne Henvisning
til en Skabelse fastsættes intet bestemt Mirakel; det overlades til
Enhver at forestille sig den dunkle Begyndelse, som han vil: Livets
Oprindelse og første Begyndelse er ingen Gjenstand for empirisk
Naturforskning.

Det gaaer hos Darwin med Læren om Skabelsen som med Læren om Midler og
Øiemed i Naturen (det Teleologiske): Skabelse og Øiemed høre ikke til
de Kategorier, der finde frugtbar Anvendelse i Erfaringslæren. Naar
Forskeren nu og da seer sig nødsaget til at anvende slige Betegnelser,
da befinder han sig i en vis Forlegenhed, fordi Ordene Skabelse,
Naturens Øiemed o.\ desl.\ maae tages paa een Gang baade i billedlig
og ubilledlig Betydning: om den egenlige Skaben som Hand-

% [Side 456]

ling, om en Tilbliven som »en Begyndelse til en Væren efter Ikkeværen«,
have vi, siger A.\ Humboldt, hverken Begreb eller Erfaring. Det Samme
maa siges om den »planmæssige og dog naturnødvendige Virksomhed, som
af Darwin tillægges den naturlige Udvælgelse«. At der endnu hviler en
mystisk Uklarhed over saadanne Grundbegreber som Skabelse og Øiemed,
det er aabenbart Filosofiens Skyld og kan ikke falde den empiriske
Naturlære til Last. Dette fører os til Betragtningen af det tredie
Hovedpunkt.

% [Side 457] (3rd main question)

\medskip

\noindent 3) Kan Darwinismen med Grund siges at indeholde
aandsfornegtende, for Religion og Sædelighed farlige Tendenser?
Dersom det kunde forudsættes, at de Mange, som i Nutiden have et
imponerende Indtryk af Naturlærens store, omfattende Betydning for
Kulturlivet, ogsaa havde en klar Erkjendelse af, hvad Naturlæren er og
maa være som Naturvidenskab, da vilde det være en let Sag at bevise,
at Darwinslæren hverken er mere eller mindre aandsfornegtende end en
hvilkensomhelst anden naturhistorisk Theori.

Naturlæren er hverken religiøs eller irreligiøs, hverken ethisk eller
uethisk; vil man fra et religiøst-ethisk Udgangspunkt angribe
Darwinismen, da angriber man med det Samme den hele vidtforgrenede
Naturvidenskab i alle dens Led og Dele.

Det skal være det Forargelige ved Darwinismen, at den vil paavise
naturlig Afstamning af høiere Arter fra lavere, fordi derved tilsigtes
en Paavisning af Menneskets naturlige Afstamning fra en ubekjendt
Dyreform. Det er, siger man, en Krænkelse af Menneskets Adel. Ganske
vist! naar man underskyder den religiøse Forestilling om en
overnaturlig Skabelsesakt og vil have Troens Forudsætninger
konfirmerede i Videnskaben. En saadan Konfirmation opnaaer man
rigtignok ikke hos Darwin; men, saa sige man os, hos hvilken
Naturforsker man da mener at opnaae den.

Louis Agassiz, Darwins erklærede Modstander, betegner rent ud hans
Afstamningslære som et videnskabeligt Misgreb, som usand i sine
Kjendsgjerninger, uvidenskabelig i sin Methode og fordærvelig i sin
Tendens: Agassiz holder strængt paa Arternes og Slægternes uforanderlige
Grændser; mon hans Theori da maanske skulde være mere religiøs end
Darwins? Agassiz gaaer ud fra Dyregeografien og udleder deraf
Konsekvenser for Menneskets Tilblivelse. Han taler om Skabelser og
Skabelsescentra og faaer ikke mindre end otte Skabelsescentra med otte
store zoologiske Riger: det arktiske, det mongolske, det europæiske,
det amerikanske Rige, Negerriget, Hottentotriget, det malaiske og det
australiske Rige. Hvert af disse Riger har sin ved en Lokalkraft
frembragte karakteristiske Fauna og sin ligeledes ved Lokalkraften
frembragte karakteristiske Menneskeæt. Her er da Skabelser nok. Langt
fra at komme i tvetydig Berøring med dyriske Overgangsformer, staaer
Mennesket selvstændig overfor Dyret; thi Mennesket er skabt, ikke blot
een Gang, men otte Gange, det kan forslaae Noget. Men at den otte Gange
gjentagne Skabelse, der skyldes den paa otte Steder værende »Lokalkraft«,
ikke har Noget med den religiøse Tro paa Skabelsen, Menneskets Skabelse
i Guds Billede, at gjøre, er øiensynligt: Agassiz's Theori er lige saa
lidt religiøs som Darwins.

% [Side 458]

Det staaer fast, at Naturvidenskaben i det Hele hverken er religiøs
eller irreligiøs, hverken hedensk eller kristelig, hverken ethisk eller
uethisk; den er, hvad den alene kan og skal være: objektivt forskende
Videnskab.

Der klages, og ikke uden Grund, over, at Videnskabens Dyrkere kun
altfor ofte misbruge Forskningens Resultater til Angreb paa Religionens
og dermed paa Livets høieste, hellige Sandheder; vi svare: naar saa
er, saa er det ikke Videnskaben, men Videnskabsdyrkerne, man har at
holde sig til.

Og hvad er saa det, der frister Videnskabsdyrkerne, de være
Darwinianere eller Antidarwinianere, til Overgreb? Det er jo netop den
gamle indgroede Fordom, at Religionens overnaturlige Kjendsgjerninger
skulde lade sig begrunde i Videnskaben, at der virkelig skulde gives en
paa Mirakler bygget Videnskab, en sand Mirakelvidenskab, som randsager
Natur og Historie, Jordens Dannelse, de levende Slægters Fremkomst,
Menneskets Oprindelse o.\ s.\ v.\ med et ganske andet dybt
videnskabeligt Blik ved Hjælp af Mirakler, end den almindelige
Forskning formaaer det uden Mirakler. Hver Gang nu den foregivne
Mirakelvidenskab sætter sin videnskabelige Mine op og vil til at vise
Forskerne til Rette, bliver dens Hulhed blottet, og idet den skarpt
tilbagevises, skeer det desværre ikke sjældent, at den høiere Sandhed,
hvis Forsvar føres med falske Vaaben, bliver miskjendt. Saalænge man
haardnakket holder paa den gamle Fordring, at Videnskaben skal være
religiøs og Religionen videnskabelig, vil det gaae stadig fremad med
Bortvendelsen fra alt Overnaturligt og Fornegtelsen af Aabenbaringens
Ord, thi Forskningens Tankegang er Evangeliets saa modsat som muligt.

Først naar det tilfulde indsees, at Religionen har og maa have sin
Erkjendelseskreds for sig selv, uafhængig af Videnskaben, og
Videnskaben sine Undersøgelseskredse for

% [Side 459]

sig, uafhængig af Religionen; først naar man ad dyrt kjøbte Erfaringers
Vei omsider har lært at indsee, at det er lige saa aandsforvirrende at
ville drage religiøse Slutninger af videnskabelige Præmisser som
videnskabelige Slutninger af religiøse Præmisser, kan der være Tale om
Fred og Forlig mellem Religion og Videnskab. Den eneste begrundede
Indvending, der kan gjøres mod en foregiven videnskabelig Theori, er,
at den er uvidenskabelig. Har Darwinsmethoden videnskabelige Mangler,
vil den blive korrigeret i Videnskaben, og Naturforskerne ere de eneste
rette, kompetente Dommere. Om den seirer eller taber i Striden, har
ingen væsenlig Indflydelse paa Aandslæren.

Det religiøse Begreb om Skabelse er ikke et Vidensbegreb, men et
Troesbegreb. Darwinismen gjør ikke Religionen Afbræk, bevæger sig ikke
paa Troens Enemærker, løser ikke Skabelsens Problem, afkræfter ikke
det aabenbarede Ord: at Mennesket er skabt i Guds Billede.

\medskip
\begin{flushright}
R.\ Nielsen
\end{flushright}

\end{document}
